%% ma: L10-B14-0029
%% lop: 10
%% chuong: 5
%% bai: 14
%% dang: 10.14.5
%% loai: DS
%% muc: [NB, NB, TH, VD]
%% dap_an: "SĐSĐ"
%% nguon: "chuyen-de-cac-so-dac-trung-cua-mau-so-lieu-khong-ghep-nhom-toan-10-chuong-trinh-moi.pdf (GV. Phan Nhật Linh) (P04), Bài 03 Dạng 1 Phần II Câu 6"
%% kien_thuc: "Lớp 10 Bài 13 (số trung bình, tứ phân vị), Bài 14 (khoảng biến thiên, khoảng tứ phân vị, giá trị bất thường)"
%% trang_thai: chinh-thuc
%% ghi_chu: "Tự duyệt theo yêu cầu giáo viên 22/10/2026."
\begin{ex}
Tiến hành đo huyết áp của $8$ người, ta thu được kết quả sau:
\begin{center}
\begin{tabular}{|c|c|c|c|c|c|c|c|}\hline
$77$ & $105$ & $117$ & $84$ & $96$ & $72$ & $105$ & $124$\\\hline
\end{tabular}
\end{center}
Xét tính đúng sai của các khẳng định sau:
\choiceTF{Khoảng biến thiên của mẫu số liệu là $50$}{\True Số trung bình của mẫu số liệu là $97{,}5$}{Khoảng tứ phân vị của mẫu số liệu là $\Delta_Q=30$}{\True Mẫu số liệu không có giá trị bất thường}
\loigiai{Sắp xếp: $72$; $77$; $84$; $96$; $105$; $105$; $117$; $124$.

a) Sai: $R=124-72=52$.

b) Đúng: $\overline x=\dfrac{780}8=97{,}5$.

c) Sai: $Q_1=\dfrac{77+84}2=80{,}5$; $Q_3=\dfrac{105+117}2=111$; $\Delta_Q=30{,}5$.

d) Đúng: $Q_1-1{,}5\Delta_Q=34{,}75$; $Q_3+1{,}5\Delta_Q=156{,}75$; mọi giá trị đều nằm trong đoạn $[34{,}75;156{,}75]$.}
\end{ex}
